\section{Get source data}

\textbf{Today.} The sources differ by workbook.
Infrastructure data come from databases: ENTSOG and GIE for Europe, the EIA for the United States and the Global Energy Monitor elsewhere, mostly as downloadable tables.
The projections are anchored to the aggregates of a published outlook, and split across nodes with national or regional statistics.
For elasticities and sector and seasonal shares, there is no data set to fetch; estimates come from country studies of varying age, or from earlier GGM work kept in what Lukas Barner called a large number of non-standardised spreadsheets.
Part of the work is already supported by tools: one developer has a script that extracts pipeline data, and has used AI tools mostly for building the lookup tables that map sources to nodes.

\textbf{What the step needs.} Finding the current edition of each source, reading it in whatever form it comes, and judging how much to trust each value.
The last of these is the hardest.
A capacity read from an operator's current database is a different kind of value from a regional consumption share inferred from GDP, and for some regions, such as parts of Russia, the developers describe the split as largely their own estimate.
Some traps require domain knowledge rather than reading skill: the documentation warns, for instance, that regional shares in the United States must be computed from marketed rather than gross production \citep[p.~24]{egging2019global}.

\textbf{What an agent could do.} Retrieval and extraction suit an LLM well (Section 4.2).
Where a source is a structured download, ordinary code should fetch and parse it, since the format is stable and the result must be exact.
Where a source is a report, a regulator's web page or a paper, an agent could find the relevant table, extract the values and record where each came from.
Lukas Barner's requirement for such an agent was specific: every value should carry its source and a statement of how it was obtained, whether from a recent official source, an older report, an analogy with a similar country (and which similarity), or a guess.
This turns the agent's output into something a reviewer can check selectively, starting with the guesses, rather than a table of numbers of unknown origin.
It also follows the design of AutoSAM, which records every value it had to assume (Section 2.2).
Licensed sources raise a separate question: if their terms do not allow sending them to an external provider, the extraction must run on a local model, as Section 4.4 discussed.

\textbf{Where a person stays involved.} Choosing which outlook anchors the scenario is part of the storyline and stays with the analysts.
Beyond that, the person's role is to review values by their confidence class rather than one by one.
The risk to watch is the plausible gap-filler: a value with no source that looks like any other value in the table, which is the hallucination problem of Section 4.3 in its most damaging form.
